
\section{Conclusion}

Our work presents a promising framework and reveals the connection between learning-based locomotion control and classical Internal Model Control principles.
In addition, our experimental results show that joint encoders and an IMU provide sufficient information.
We also demonstrate that compared to regression, contrastive learning provides a better representation in terms of robustness to noise and high sample efficiency.
To the best of our knowledge, our method is the first system to fully utilize batch-level information in locomotion tasks, leveraging massively parallel simulations of Isaac Gym.
However, the lack of external perceptions, e.g., visual information, leaves legged robots with only a primitive skill set.
We will incorporate our method with external perception in the future to overcome more challenging environments and higher-level tasks.

\paragraph{Acknowledgement.} 
This work is supported by Shanghai Artificial Intelligence Laboratory.